% ============================================================
% 交流回路（オシロスコープを用いた測定） 記録用ワークシート本文
% 原本: 実験手引書_交流回路(オシロスコープを用いた測定).docx の末尾の記録表
% ============================================================

\begin{center}
  {\sffamily\gtfamily\bfseries\Large 電気電子工学実験　交流回路（オシロスコープ）記録用紙}
\end{center}

\begin{flushright}
  学籍番号：\underline{\hspace{12\zw}}\quad
  氏名：\underline{\hspace{12\zw}}
\end{flushright}

% TODO(原文確認): 手順では 20 Hz～1 kHz だが，表は 1 Hz から
\begin{table}[h]
  \centering
  \caption{実験1：$RC$ローパスフィルタの特性}
  \label{tab:oscilloscope-rc}
  \renewcommand{\arraystretch}{1.45}
  \begin{tabular}{|c|>{\centering\arraybackslash}p{8\zw}|>{\centering\arraybackslash}p{8\zw}|>{\centering\arraybackslash}p{8\zw}|}
    \hline
    周波数 [Hz] & 入力電圧 $V_\mathrm{in}$ [V] & 出力電圧 $V_\mathrm{out}$ [V] & 入出力比 $V_\mathrm{out}/V_\mathrm{in}$ \\ \hline\hline
    1 & & & \\ \hline
    10 & & & \\ \hline
    100 & & & \\ \hline
    200 & & & \\ \hline
    300 & & & \\ \hline
    400 & & & \\ \hline
    500 & & & \\ \hline
    600 & & & \\ \hline
    700 & & & \\ \hline
    800 & & & \\ \hline
    900 & & & \\ \hline
    1000 & & & \\ \hline
  \end{tabular}
\end{table}

\clearpage

\begin{table}[h]
  \centering
  \caption{実験2：$RLC$直列共振フィルタの特性}
  \label{tab:oscilloscope-rlc}
  \renewcommand{\arraystretch}{1.45}
  \begin{tabular}{|c|>{\centering\arraybackslash}p{8\zw}|>{\centering\arraybackslash}p{8\zw}|>{\centering\arraybackslash}p{8\zw}|}
    \hline
    周波数 [Hz] & 入力電圧 $V_\mathrm{in}$ [V] & 出力電圧 $V_\mathrm{out}$ [V] & 入出力比 $V_\mathrm{out}/V_\mathrm{in}$ \\ \hline\hline
    2000 & & & \\ \hline
    2500 & & & \\ \hline
    3000 & & & \\ \hline
    3500 & & & \\ \hline
    4000 & & & \\ \hline
    4500 & & & \\ \hline
    5000 & & & \\ \hline
    5500 & & & \\ \hline
    6000 & & & \\ \hline
    6500 & & & \\ \hline
    7000 & & & \\ \hline
    7500 & & & \\ \hline
    8000 & & & \\ \hline
  \end{tabular}
\end{table}
